\documentclass[11pt]{article}

\usepackage[margin=0.82in]{geometry}
\usepackage{booktabs}
\usepackage{tabularx}
\usepackage{array}
\usepackage{amsmath}
\usepackage{graphicx}
\usepackage{xcolor}
\usepackage{enumitem}
\usepackage{microtype}
\usepackage{fancyhdr}
\usepackage{hyperref}

\hypersetup{
    colorlinks=true,
    linkcolor=blue!55!black,
    urlcolor=blue!60!black,
    citecolor=blue!55!black
}

\pagestyle{fancy}
\fancyhf{}
\lhead{\small First-Stage Research Note}
\rhead{\small Bob Shen}
\cfoot{\thepage}
\renewcommand{\headrulewidth}{0.3pt}

\setlength{\parindent}{0pt}
\setlength{\parskip}{0.55em}
\setlist[itemize]{leftmargin=1.35em, itemsep=0.25em, topsep=0.25em}
\setlist[enumerate]{leftmargin=1.5em, itemsep=0.25em, topsep=0.25em}

\newcommand{\realdata}{\textcolor{green!40!black}{\textbf{Real data}}}
\newcommand{\synthetic}{\textcolor{orange!75!black}{\textbf{Synthetic only}}}
\newcommand{\pending}{\textcolor{red!65!black}{\textbf{Pending}}}

\title{
    \textbf{First-Stage Research Note}\\[0.25em]
    \large Thermal Extremes and Cardiovascular Admissions in Hong Kong,
    2013--2023
}
\author{
    Bob Shen Ruililin\\
    Laidlaw Scholars Programme, The University of Hong Kong
}
\date{10 July 2026}

\begin{document}

\maketitle

\begin{center}
\small
\textit{
This note distinguishes real public-data findings from synthetic workflow
testing. No real Hospital Authority outcome association has yet been estimated.
}
\end{center}

\section*{Executive Summary}

During the first stage of the project, I built and executed a reproducible
R workflow intended to support the eventual analysis of monthly Hospital
Authority (HA) cardiovascular admissions. I also processed real Hong Kong
Observatory (HKO) weather data into monthly averages and officially defined
thermal-extreme measures.

The principal findings at this stage are methodological and descriptive:

\begin{itemize}
    \item The HKO exposure-processing pipeline reproduces published annual
    counts of hot nights, very hot days, and cold days for all 11 years
    2013--2023 (33/33 annual totals match HKO \textit{Year's Weather} summaries).
    \item The study period contains substantial variation in thermal extremes.
    For example, HKO recorded 10 hot nights and 17 very hot days in 2013,
    compared with 61 hot nights and 54 very hot days in 2021.
    \item Cold exposure has not disappeared. HKO recorded 14 cold days in
    2013 and 13 in 2021. The appropriate scientific question is therefore not
    whether heat has simply replaced cold, but whether an emerging heat burden
    now coexists with persistent cold-related risk.
    \item The planned health-outcome workflow has been tested using synthetic
    counts. This validates the mechanics of data merging, offsets, count
    models, diagnostics, and output generation, but provides no evidence about
    actual AMI or stroke admissions.
    \item The primary remaining dependency is a clearly defined, securely
    accessible HA outcome extract. Public C\&SD population and EPD pollution
    data can be completed in parallel.
\end{itemize}

The next scientific milestone is therefore not a more elaborate model. It is
agreement on the HA outcome definition, episode structure, age bands, and
secure data-access process.

\section{Research Question and Current Framing}

Earlier Hong Kong research used daily observations and reported a strong
association between cooler temperatures and acute myocardial infarction
(AMI) admissions. Goggins et al.\ found that, below an estimated threshold of
$24^\circ$C, each $1^\circ$C decrease was associated with a 3.7\% increase in
Hong Kong AMI admissions using an average temperature lag of 0--13 days.
No statistically significant heat effect was detected, and same-day
nitrogen dioxide was the strongest pollutant predictor \cite{goggins2013}.

The present project asks whether the contemporary cardiovascular risk profile
remains predominantly cold-related or now includes detectable associations
with increasingly frequent heat extremes. The proposed study population is
Hong Kong residents aged 35 years and older, and the intended outcomes are
monthly admissions for AMI, ischemic stroke, and hemorrhagic stroke.

The most defensible current framing is:

\begin{quote}
\textit{How are officially defined thermal extremes---particularly hot nights
and cold days---associated with contemporary cardiovascular hospital burden
in an aging Hong Kong, and how do these associations vary by diagnosis and
age?}
\end{quote}

This framing does not assume in advance that a ``regime shift'' has occurred.
A persistent cold-dominant association would remain scientifically and
operationally important.

\section{Work Completed}

\subsection{Climate-data processing}

I built the weather-processing portion of the R pipeline using HKO daily
observations at the Observatory Headquarters. It converts daily data into
monthly variables including mean, minimum, and maximum temperature; relative
humidity and rainfall; hot nights ($T_{\min}\ge 28^\circ$C); very hot days
($T_{\max}\ge 33^\circ$C); extremely hot days ($T_{\max}\ge 35^\circ$C); cold
days ($T_{\min}\le 12^\circ$C); and longest consecutive runs of hot nights and
very hot days.

The pipeline includes date coverage, missingness, plausible-value, and
annual-total checks. Annual extreme counts for 2013--2023 agree with the
corresponding published HKO summaries (see Figure~\ref{fig:hko-extremes} and
the validation table in the repository).

\subsection{Statistical dry run}

I created an explicitly synthetic HA-like panel with structure
$\mathrm{month}\times\mathrm{age}\times\mathrm{sex}\times\mathrm{diagnosis}$
to test merging, population$\times$days offsets, Poisson/Negative Binomial
specifications, interactions, COVID sensitivities, diagnostics, and automatic
table/figure generation. This demonstrates that the workflow can accept an HA
file once its structure is finalized. It does \textbf{not} produce substantive
cardiovascular findings.

\subsection{Research documents}

The repository also contains draft materials for the analysis plan and
assumption log; data requests for Roro and Hogan; a preliminary Methods
section; table and figure shells; and a literature baseline. The purpose of
the present note is to summarize the scientifically relevant content without
requiring the reader to inspect the code.

\section{Preliminary Findings from Real HKO Data}

Table~\ref{tab:hko} presents selected benchmark years from official HKO annual
summaries. It is not intended as a formal trend analysis.

\begin{table}[ht]
\centering
\caption{Selected HKO thermal-extreme benchmarks}
\label{tab:hko}
\small
\begin{tabular}{lrrrrl}
\toprule
Year & Mean temp. & Hot nights & Very hot days & Cold days & Status\\
     & ($^\circ$C) & & & &\\
\midrule
2013 & 23.3 & 10 & 17 & 14 & Core period\\
2017 & 23.9 & 41 & 29 & 9  & Core period\\
2021 & 24.6 & 61 & 54 & 13 & Core period\\
2024 & 24.8 & 50 & 52 & 11 & Context only\\
\bottomrule
\end{tabular}
\end{table}

\subsection*{1. The study period contains meaningful heat-exposure contrast}

The later study period contains substantially more hot nights and very hot
days than 2013. This is useful analytically: the 2013--2023 period is not
uniformly hot and therefore contains exposure variation with which to estimate
associations. This should not yet be described as a formal trend finding.

\subsection*{2. Hot nights are a defensible primary heat metric}

Monthly average temperature can obscure short episodes of extreme exposure.
Hot nights measure the absence of overnight thermal relief and are defined
consistently by HKO. Very hot days, mean minimum temperature, heat-spell
duration, and wet-bulb temperature should be considered alternative exposure
specifications rather than entered simultaneously into one highly collinear
primary model.

\subsection*{3. Cold remains part of the exposure distribution}

Cold days remained present even in warm years. The project should therefore
retain both heat and cold metrics and avoid framing the analysis as a forced
competition in which one must replace the other.

\begin{figure}[htbp]
    \centering
    \includegraphics[width=0.91\textwidth]
    {figures/hko_annual_extremes_2013_2023.pdf}
    \caption{
    Annual counts of hot nights, very hot days, and cold days at the
    Hong Kong Observatory Headquarters, 2013--2023. Hot nights are defined as days
    with a minimum temperature of at least $28^\circ$C; very hot days
    as days with a maximum temperature of at least $33^\circ$C; and
    cold days as days with a minimum temperature of at most
    $12^\circ$C. These are environmental exposure descriptors and not
    health-outcome results. Pipeline annual totals match official HKO
    \textit{Year's Weather} summaries for all years shown (33/33 checks).
    }
    \label{fig:hko-extremes}
\end{figure}

\section{Lessons from the Coding and Dry Run}

\subsection{Population offsets are essential}

Because the older population changed substantially during the study period,
raw admission counts cannot be interpreted as changing risk. The proposed
offset is $\log(P_{ast}\times D_t)$, where $P_{ast}$ is the population for
age group $a$, sex $s$, and month $t$, and $D_t$ is days in that month.

\subsection{The effective environmental sample is 132 months}

Although stratification may produce several thousand outcome rows, the
territory-wide environmental exposures vary across only 132 months. Inference
will need to recognize common monthly exposure and repeated observations
within outcome strata.

\subsection{Monthly data favor burden estimation over acute-trigger analysis}

Earlier Hong Kong studies used daily data and distributed lag structures.
Heat effects may occur over short lags, while cold effects may persist over
longer periods. Monthly aggregation may therefore attenuate heat associations
more strongly than cold associations. Results should be interpreted as monthly
hospital-burden associations, not estimates of same-day physiological
triggering.

\subsection{Outcome definition matters as much as model choice}

Before analysis, it is necessary to determine whether the file counts
admissions, patients, or completed episodes; whether principal diagnosis is
available; how transfers and readmissions are handled; whether ED attendance
followed by admission is counted twice; and whether emergency and elective
admissions can be distinguished.

\section{Proposed Statistical Framework}

The primary analyses should be diagnosis-specific. A provisional Negative
Binomial model includes heat exposure $H_t$ (initially monthly hot nights),
cold exposure $C_t$, age, sex, diagnosis, month effects, a smooth time trend,
prespecified covariates $\mathbf{Z}_t$, and the population--time offset.
Poisson regression should be fitted as a baseline and checked for
overdispersion. Principal effect estimates would be rate ratios per five
additional hot nights or cold days. Interactions (thermal$\times$age,
thermal$\times$diagnosis, heat$\times$ozone) should be introduced separately.

\section{Data Status and Transfer}

\begin{table}[ht]
\centering
\caption{Current data status}
\small
\begin{tabularx}{\textwidth}{
    >{\raggedright\arraybackslash}p{2.2cm}
    >{\raggedright\arraybackslash}p{2.0cm}
    X
    >{\raggedright\arraybackslash}p{2.6cm}}
\toprule
Source & Status & What is needed & Proposed handling\\
\midrule
HKO weather &
\realdata\ processed &
Complete final quality review; add wet-bulb or multi-station measures only if
they improve the prespecified analysis. &
Public data; institutional email, repository, or shared drive.\\
C\&SD population &
\pending &
Download the half-yearly age-sex population table and construct monthly
denominators, retaining 65--69 and 70--74 where possible. &
Public CSV/XLSX/JSON; shared drive or repository.\\
EPD pollution &
\pending &
Download NO$_2$, O$_3$, PM$_{2.5}$, and PM$_{10}$ for selected general
stations (monthly portal limit implies $\ge$2 exports for 132 months). &
Public CSV; shared drive or institutional email.\\
HA outcomes &
\pending &
Confirm coding, principal diagnosis, episode structure, age bands, suppression,
and approved analysis environment. &
Approved HA/HKU secure process only unless Roro confirms otherwise.\\
Synthetic outcomes &
\synthetic &
Retain only for pipeline tests; never merge into substantive output. &
Repository, clearly labelled synthetic.\\
\bottomrule
\end{tabularx}
\end{table}

Public environmental and population files can ordinarily be exchanged through
institutional email or a shared drive. No patient-level, restricted,
small-cell, or unvetted HA output should be sent by ordinary email until Roro
confirms governance status.

\section{Questions That Require Team Input}

Four issues are immediate blockers:

\begin{enumerate}
    \item \textbf{Coding and diagnosis position.} ICD-9-CM, ICD-10, or local
    coding? Can principal diagnoses be identified?
    \item \textbf{Unit of outcome.} Admissions, unique patients, or completed
    episodes? How are transfers and readmissions handled?
    \item \textbf{Granularity and suppression.} Can five-year age bands be
    supplied, especially 65--69 and 70--74? What small-cell rules apply?
    \item \textbf{Secure access.} Secure-server analysis versus approved
    aggregate transfer? Which packages and output-vetting procedures apply?
\end{enumerate}

Choices that can be deferred without delaying the project include extending
outcomes through 2024, district analyses, comorbidity extensions, the final
pollution station set, influenza in primary versus sensitivity models, and the
exact correlated-count model for final inference. My current default is to
retain 2013--2023 as the core period and request 2024 only if readily
available.

\section{Proposed Next Steps}

\begin{enumerate}
    \item Send Roro focused HA extraction questions and request a mock schema
    or data dictionary.
    \item Download and process real C\&SD population data.
    \item Obtain EPD files for NO$_2$, O$_3$, PM$_{2.5}$, and PM$_{10}$
    (general stations first).
    \item Produce a final environmental QC report (missingness, correlations,
    annual exposure plots).
    \item Freeze a short pre-analysis plan after the HA schema is confirmed.
    \item Run descriptive validation of HA counts before regression.
    \item Fit the prespecified model ladder.
\end{enumerate}

I have also begun a structured literature update, which I will use to refine
the Introduction and prespecify historical comparisons before examining HA
results.

\section{Proposed Meeting}

A 30-minute meeting with Dr.\ Bishai, Hogan, Roro, and Bob would be sufficient
to unlock the next stage: confirm the primary question and 2013--2023 core
period; confirm HA outcome, coding, and episode structure; confirm age bands
and suppression; confirm secure access; allocate C\&SD/EPD responsibility; and
agree a date for the first validated HA descriptive report.

\section*{Conclusion}

The first-stage work has established that the environmental exposure pipeline
is feasible, reproducible, and capable of representing both heat and cold
extremes. It has also produced an analysis framework ready to receive the
clinical and population files. The project has not yet produced a
cardiovascular association result. The next step should prioritize outcome
validity and secure data integration rather than additional model complexity.

\clearpage
\appendix
\section{Working Manuscript Introduction}

Temperature-related cardiovascular morbidity is an important concern in
subtropical cities, where populations may experience both intense humid heat
and episodic cold despite generally warm climates. Acute myocardial infarction
and stroke may respond differently to thermal exposure, and the observed
burden can vary according to age, housing, behavioral adaptation, air
pollution, and access to temperature-controlled environments. These questions
are especially relevant in Hong Kong, a densely populated and rapidly aging
city with substantial seasonal variation and a growing frequency of hot
nights and very hot days.

Earlier Hong Kong evidence emphasized cold-related cardiovascular risk.
Using daily data from 2000--2009, Goggins et al.\ reported that decreases in
temperature below an estimated $24^\circ$C threshold were associated with
increased AMI admissions, while no statistically significant heat effect was
observed \cite{goggins2013}. The same study identified nitrogen dioxide as an
important predictor of AMI admissions.

Hong Kong's thermal environment has changed since these historical study
periods. At the Hong Kong Observatory, the number of hot nights increased from
10 in 2013 to 61 in 2021, while very hot days increased from 17 to 54 in the
same selected-year comparison \cite{hko2013,hko2021}. Cold days nevertheless
continued to occur, indicating that contemporary preparedness must consider
potentially overlapping heat and cold hazards.

Demographic change may further alter the distribution of cardiovascular
burden. Age-sex-specific population denominators and appropriate person-time
offsets are necessary to distinguish changing population structure from
changes in admission rates. The environmental context has also evolved:
temperature and pollutants should not be treated as statistically
interchangeable covariates, and ozone in particular may act as a confounder,
effect modifier, or pathway variable.

This study will examine monthly AMI, ischemic stroke, and hemorrhagic stroke
admissions among Hong Kong residents aged 35 years or older from 2013 to 2023.
It will assess associations with HKO-defined thermal extremes, emphasizing hot
nights and cold days, while accounting for population denominators,
seasonality, long-term trends, and COVID-period disruption. The objective is
not to assume that heat has replaced cold, but to determine whether the
historically cold-dominant pattern remains sufficient to describe
contemporary cardiovascular hospital burden in a warming and aging Hong Kong.

\clearpage
\section{Use of AI-Assisted Tools}

AI-assisted tools were used for code scaffolding, methodological
brainstorming, and review of potential data and modeling issues. Synthetic
data were used to test the workflow before access to HA outcomes. All
substantive claims about HKO annual extremes were checked against official
HKO \textit{Year's Weather} pages. Literature citations used in this note
were restricted to verified primary sources; a fuller evidence matrix is under
construction as an internal deliverable.

\clearpage
\begin{thebibliography}{9}

\bibitem{goggins2013}
Goggins WB, Chan EYY, Yang C-Y.
Weather, pollution, and acute myocardial infarction in Hong Kong and Taiwan.
\textit{International Journal of Cardiology}.
2013;168(1):243--249.
doi:10.1016/j.ijcard.2012.09.087.

\bibitem{hko2013}
Hong Kong Observatory.
\textit{The Year's Weather -- 2013}.
\url{https://www.hko.gov.hk/en/wxinfo/pastwx/ywx2013.htm}

\bibitem{hko2017}
Hong Kong Observatory.
\textit{The Year's Weather -- 2017}.
\url{https://www.hko.gov.hk/en/wxinfo/pastwx/2017/ywx2017.htm}

\bibitem{hko2021}
Hong Kong Observatory.
\textit{The Year's Weather -- 2021}.
\url{https://www.hko.gov.hk/en/wxinfo/pastwx/2021/ywx2021.htm}

\bibitem{hko2024}
Hong Kong Observatory.
\textit{The Year's Weather -- 2024}.
\url{https://www.hko.gov.hk/en/wxinfo/pastwx/2024/ywx2024.htm}

\bibitem{hkoopendata}
Hong Kong Observatory.
\textit{Hong Kong Observatory Open Data}.
\url{https://www.hko.gov.hk/en/abouthko/opendata_intro.htm}

\bibitem{epd}
Environmental Protection Department.
\textit{Environmental Protection Interactive Centre: Air Quality Data}.
\url{https://cd.epic.epd.gov.hk/EPICDI/air/station/?lang=en}

\bibitem{csd}
Census and Statistics Department.
\textit{Hong Kong Population: Population by Sex and Age Group}.
\url{https://data.gov.hk/en-data/dataset/hk-censtatd-tablechart-110-01001}

\end{thebibliography}

\end{document}
